Consider an isotropic light source of frequency $f_R$ in a frame which is
fixed to the source (i.e.\ rest frame). In this rest frame, consider a light
ray emitted from the source that makes an angle $\theta_R$ with the $X$-axis.
The light source is moving along positive $X$ direction with (relativistic)
speed $v$ as measured in the lab frame.

\begin{description}
\item[(a)] Find an expression for the frequency $f_L$ of this ray in the lab
  frame, and the cosine of the angle that this ray makes with the $X$-axis in
  the lab frame. \hfill[11\,pt]\\
  \textit{Hint:} In relativistic mechanics, energy $E$ and momentum $p$ of a
  particle between rest and lab frame are related in the following way:
  \begin{align*}
    \frac{E_L}{c} &= \gamma\left(\frac{E_R}{c} + p_{x_R}\frac{v}{c}\right)\\
    p_{x_L} &= \gamma\left(p_{x_R} + \frac{E_R v}{c^2}\right)\\
    p_{y_L} &= p_{y_R}\\
    p_{z_L} &= p_{z_R}
  \end{align*}
  where:
  \[ \gamma = \frac{1}{\sqrt{1 - \frac{v^2}{c^2}}} \]
\item[(b)] For the following cases:
  \begin{description}
  \item[i)] $\theta_R = 0^\circ$
  \item[ii)] $\theta_R = \cos^{-1}(-v/c)$
  \item[iii)] $\theta_R = 90^\circ$
  \item[iv)] $\theta_L = 180^\circ$
  \end{description}
  draw direction vectors of the beam in $XY$ plane of the rest frame as well
  as separately in $X'Y'$ plane of the lab frame. \hfill[4\,pt]
\end{description}

In accretion disks around black holes, the charged particles are orbiting at
relativistic speeds and in their rest frames may be considered as isotropic
point sources of light. Consider such a particle K in a circular orbit of
radius $r$ and angular speed $\omega$ around a central object located at $O$
(see figure).

\ioaafig{2022_TH_13-f1.pdf}

Let us assume that our lab frame is fixed to an observer located at a point
$P$ on the $OY$ axis, which is stationary with respect to $O$.
$OP = R \gg r$. Let $t_{L0} = t_{R0} = 0$ correspond to the moment when K is
seen crossing the $OX$ axis. As K is moving with relativistic speed, the
duration $\Delta t_R$ measured by an observer in the rest frame of the source
K is related to the duration measured in the lab frame $\Delta t_L$ at $P$ by
the expression $\Delta t_L = \gamma\Delta t_R$.

\begin{description}
\item[(c)] Derive an expression for $f_L$ as a function of $t_L$
  ($t_L > R/c$)? \hfill[7\,pt]
\end{description}

Let us consider a fraction of the light from the source that is emitted in an
infinitesimal solid angle
$\Delta\Omega_R = -\Delta(\cos\theta_R) \cdot \Delta\phi$ in the direction
making an angle $\theta_R$ with respect to the $X$ axis in the rest frame, as
it is shown on the figure below.

\ioaafig{2022_TH_13-f2.pdf}

\begin{description}
\item[(d)] Show that, as measured in the lab frame
  \[ \Delta\Omega_L = \frac{\Delta\Omega_R}
     {\gamma^2\left(1 + \frac{v}{c}\cos\theta_R\right)^2} \]
  \hfill[10\,pt]
\item[(e)] If the intrinsic luminosity of the light source is $L$, what is the
  energy flux $F_L$ observed by the observer at point $P$ at the moment $t_L$
  ($t_L > R/c$)? \hfill[15\,pt]\\
  \textit{Hint:} In the rest frame of the source, you may assume $N_R$ number
  of photons are directed within the solid angle $\Delta\Omega_R$ during the
  time interval $\Delta t_R$.
\item[(f)] Charged particles in the relativistic beam shot from the
  supermassive black hole at the centre of the galaxy M87 have speeds up to
  $0.95c$. What would be the maximum and minimum amplification factor for the
  energy flux for a relativistic beam from M87? \hfill[3\,pt]
\end{description}
